% Options for packages loaded elsewhere
\PassOptionsToPackage{unicode}{hyperref}
\PassOptionsToPackage{hyphens}{url}
\documentclass[
  11pt,
]{article}
\usepackage{xcolor}
\usepackage[margin=1in]{geometry}
\usepackage{amsmath,amssymb}
\setcounter{secnumdepth}{-\maxdimen} % remove section numbering
\usepackage{iftex}
\ifPDFTeX
  \usepackage[T1]{fontenc}
  \usepackage[utf8]{inputenc}
  \usepackage{textcomp} % provide euro and other symbols
\else % if luatex or xetex
  \usepackage{unicode-math} % this also loads fontspec
  \defaultfontfeatures{Scale=MatchLowercase}
  \defaultfontfeatures[\rmfamily]{Ligatures=TeX,Scale=1}
\fi
\usepackage{lmodern}
\ifPDFTeX\else
  % xetex/luatex font selection
\fi
% Use upquote if available, for straight quotes in verbatim environments
\IfFileExists{upquote.sty}{\usepackage{upquote}}{}
\IfFileExists{microtype.sty}{% use microtype if available
  \usepackage[]{microtype}
  \UseMicrotypeSet[protrusion]{basicmath} % disable protrusion for tt fonts
}{}
\makeatletter
\@ifundefined{KOMAClassName}{% if non-KOMA class
  \IfFileExists{parskip.sty}{%
    \usepackage{parskip}
  }{% else
    \setlength{\parindent}{0pt}
    \setlength{\parskip}{6pt plus 2pt minus 1pt}}
}{% if KOMA class
  \KOMAoptions{parskip=half}}
\makeatother
\usepackage{color}
\usepackage{fancyvrb}
\newcommand{\VerbBar}{|}
\newcommand{\VERB}{\Verb[commandchars=\\\{\}]}
\DefineVerbatimEnvironment{Highlighting}{Verbatim}{commandchars=\\\{\}}
% Add ',fontsize=\small' for more characters per line
\newenvironment{Shaded}{}{}
\newcommand{\AlertTok}[1]{\textcolor[rgb]{1.00,0.00,0.00}{\textbf{#1}}}
\newcommand{\AnnotationTok}[1]{\textcolor[rgb]{0.38,0.63,0.69}{\textbf{\textit{#1}}}}
\newcommand{\AttributeTok}[1]{\textcolor[rgb]{0.49,0.56,0.16}{#1}}
\newcommand{\BaseNTok}[1]{\textcolor[rgb]{0.25,0.63,0.44}{#1}}
\newcommand{\BuiltInTok}[1]{\textcolor[rgb]{0.00,0.50,0.00}{#1}}
\newcommand{\CharTok}[1]{\textcolor[rgb]{0.25,0.44,0.63}{#1}}
\newcommand{\CommentTok}[1]{\textcolor[rgb]{0.38,0.63,0.69}{\textit{#1}}}
\newcommand{\CommentVarTok}[1]{\textcolor[rgb]{0.38,0.63,0.69}{\textbf{\textit{#1}}}}
\newcommand{\ConstantTok}[1]{\textcolor[rgb]{0.53,0.00,0.00}{#1}}
\newcommand{\ControlFlowTok}[1]{\textcolor[rgb]{0.00,0.44,0.13}{\textbf{#1}}}
\newcommand{\DataTypeTok}[1]{\textcolor[rgb]{0.56,0.13,0.00}{#1}}
\newcommand{\DecValTok}[1]{\textcolor[rgb]{0.25,0.63,0.44}{#1}}
\newcommand{\DocumentationTok}[1]{\textcolor[rgb]{0.73,0.13,0.13}{\textit{#1}}}
\newcommand{\ErrorTok}[1]{\textcolor[rgb]{1.00,0.00,0.00}{\textbf{#1}}}
\newcommand{\ExtensionTok}[1]{#1}
\newcommand{\FloatTok}[1]{\textcolor[rgb]{0.25,0.63,0.44}{#1}}
\newcommand{\FunctionTok}[1]{\textcolor[rgb]{0.02,0.16,0.49}{#1}}
\newcommand{\ImportTok}[1]{\textcolor[rgb]{0.00,0.50,0.00}{\textbf{#1}}}
\newcommand{\InformationTok}[1]{\textcolor[rgb]{0.38,0.63,0.69}{\textbf{\textit{#1}}}}
\newcommand{\KeywordTok}[1]{\textcolor[rgb]{0.00,0.44,0.13}{\textbf{#1}}}
\newcommand{\NormalTok}[1]{#1}
\newcommand{\OperatorTok}[1]{\textcolor[rgb]{0.40,0.40,0.40}{#1}}
\newcommand{\OtherTok}[1]{\textcolor[rgb]{0.00,0.44,0.13}{#1}}
\newcommand{\PreprocessorTok}[1]{\textcolor[rgb]{0.74,0.48,0.00}{#1}}
\newcommand{\RegionMarkerTok}[1]{#1}
\newcommand{\SpecialCharTok}[1]{\textcolor[rgb]{0.25,0.44,0.63}{#1}}
\newcommand{\SpecialStringTok}[1]{\textcolor[rgb]{0.73,0.40,0.53}{#1}}
\newcommand{\StringTok}[1]{\textcolor[rgb]{0.25,0.44,0.63}{#1}}
\newcommand{\VariableTok}[1]{\textcolor[rgb]{0.10,0.09,0.49}{#1}}
\newcommand{\VerbatimStringTok}[1]{\textcolor[rgb]{0.25,0.44,0.63}{#1}}
\newcommand{\WarningTok}[1]{\textcolor[rgb]{0.38,0.63,0.69}{\textbf{\textit{#1}}}}
\usepackage{longtable,booktabs,array}
\usepackage{calc} % for calculating minipage widths
% Correct order of tables after \paragraph or \subparagraph
\usepackage{etoolbox}
\makeatletter
\patchcmd\longtable{\par}{\if@noskipsec\mbox{}\fi\par}{}{}
\makeatother
% Allow footnotes in longtable head/foot
\IfFileExists{footnotehyper.sty}{\usepackage{footnotehyper}}{\usepackage{footnote}}
\makesavenoteenv{longtable}
\setlength{\emergencystretch}{3em} % prevent overfull lines
\providecommand{\tightlist}{%
  \setlength{\itemsep}{0pt}\setlength{\parskip}{0pt}}
\usepackage{bookmark}
\IfFileExists{xurl.sty}{\usepackage{xurl}}{} % add URL line breaks if available
\urlstyle{same}
\hypersetup{
  hidelinks,
  pdfcreator={LaTeX via pandoc}}

\title{Topology Determines Emergent Self-Prediction in Multi-Agent Communication Networks: A Scaling Law for Consciousness Correlate Collapse}
\author{Brian Munene Mwirigi}
\date{}

\begin{document}
\maketitle

\begin{center}\rule{0.5\linewidth}{0.5pt}\end{center}

\subsection{Abstract}\label{abstract}

This paper presents \emph{consim}, a multi-agent simulation that studies
how self-prediction emerges from lossy predictive communication on graph
topologies. The system is inspired by the Machine Consciousness
Hypothesis (Fitz, 2025). It places distributed agents on a network where
they communicate through noisy channels and learn to predict their
neighbors' states, but never their own. I measure four
consciousness-relevant quantities from the MCH framework: integration
(\(\Phi\)), reflexivity (\(R\)), temporal persistence (\(T\)), and
causal efficacy (\(E\)). Across 4,180 simulation runs on five
topologies, five noise levels, five grid sizes (12, 18, 24, 36, 48), and
three activation functions (tanh, linear, ReLU), I report three main
results. First, on high-connectivity grids, causal efficacy
anti-correlates with self-prediction (\(r = -0.71\)) and with temporal
persistence (\(r = -0.82\)). Self-determination destroys self-knowledge,
which inverts assumptions in IIT (Tononi, 2004), MCH (Fitz, 2025), and
autonomy-based consciousness frameworks. This trade-off holds across all
three activations (tanh \(-0.82\), linear \(-0.77\), ReLU \(-0.76\)) and
persists identically in a null model without learning (\(r = -0.79\)),
making it a topological invariant rather than a learning artifact.
Second, PCA on the five metrics shows they collapse to a single
principal component on Moore grids (82.5\% variance explained), forming
an ``autonomy-predictability axis.'' Third, this collapse follows a
scaling law: \(\text{PC1} = -0.136 + 0.124 \times \ln(N)\)
(\(R^2 = 0.987\)) across five grid sizes (144 to 2,304 agents), with a
sharp inflection between \(N = 324\) and \(N = 576\) agents. Random
topologies show no collapse at any size. I present a formal argument
that this works like an uncertainty principle for networked systems: on
\(K\)-regular graphs with bounded activations and correlated neighbors,
variance reduction through neighbor-averaging forces a trade-off between
causal autonomy and predictability. The constraint on consciousness
correlates is an emergent property of structured networks that only
appears above a critical system size. Code and data:
github.com/brian-mwirigi/consim.

\textbf{Keywords:} consciousness correlates, integrated information,
multi-agent systems, emergent self-prediction, network topology, scaling
laws, dimensionality collapse

\begin{center}\rule{0.5\linewidth}{0.5pt}\end{center}

\subsection{1. Introduction}\label{introduction}

Theories of consciousness use multiple measurable correlates:
integration (Tononi, 2004; Tononi et al., 2016), temporal stability
(Baars, 1988), causal efficacy (Seth, 2014; Seth \& Bayne, 2022), and
self-modeling (Graziano, 2013; Cleeremans, 2011). These theories tend to
assume the correlates can be independently optimized. The Machine
Consciousness Hypothesis (MCH) proposes that consciousness is a
substrate-free functional property of computational systems capable of
second-order perception. It arises when groups of local observers
exchange lossy predictive messages in a universal self-organizing
environment (Fitz, 2025). The MCH defines four measurable correlates:
integration (\(\Phi\)), reflexivity (\(R\)), temporal persistence
(\(T\)), and causal efficacy (\(E\)).

The question I wanted to answer is whether these correlates are
genuinely independent. If the network structure that enables them also
constrains their joint achievability, then counting independent
correlates overcounts the degrees of freedom available to any networked
system.

I address this computationally using \emph{consim}, a multi-agent
simulation that places \(N = \text{size}^2\) agents on a toroidal grid.
Each agent has an 8-dimensional internal state and a learned weight
matrix. Every tick, agents broadcast transformed state messages, receive
noisy versions of their neighbors' messages, update their states, and
learn via gradient descent on neighbor prediction error. Self-prediction
(the cosine similarity between an agent's broadcast and its next state)
is never explicitly trained. If it rises above chance, the agent has
built an implicit model of its own dynamics as a side effect of modeling
its neighbors. This is how I operationalize Fitz's idea of reflexivity
emerging from predictive communication.

Previous work on integrated information and phase transitions has
focused on Ising-model systems (Khajehabdollahi, 2018; Luitle, 2023) or
coupled oscillators (Abrego \& Zaikin, 2019), looking at how \(\Phi\)
alone varies with temperature or coupling strength. Related theoretical
frameworks propose morphospaces for consciousness (Arsiwalla et al.,
2023) or spectral measures of emergence (Bailey \& Schneider, 2025). The
MCH itself (Fitz, 2025) is a purely theoretical proposal with no
empirical implementation.

My contribution is different from prior work in four ways: (1) I measure
five consciousness correlates at the same time, not just \(\Phi\); (2) I
vary network topology rather than temperature, using structured and
random graphs; (3) I show that these correlates are not independent on
high-connectivity networks, collapsing to a single principal component;
and (4) I find a scaling law for this collapse, with a phase-like
transition between \(N = 324\) and \(N = 576\) agents on structured
grids. No previous work has shown dimensionality collapse of
consciousness correlates, the T-E autonomy-predictability trade-off, or
the functional identity E \(\equiv\) \(\Phi\) on structured networks.

\begin{center}\rule{0.5\linewidth}{0.5pt}\end{center}

\subsection{2. Methods}\label{methods}

\subsubsection{2.1 Agent Architecture}\label{agent-architecture}

Each agent \(i\) has state \(s_i \in \mathbb{R}^8\) and weights
\(W_i \in \mathbb{R}^{8 \times 8}\), initialized randomly. Every tick,
these operations happen:

\begin{enumerate}
\def\labelenumi{\arabic{enumi}.}
\tightlist
\item
  \textbf{Broadcast:} \(m_i = \sigma(W_i \cdot s_i)\) where \(\sigma\)
  is a configurable activation function
\item
  \textbf{Noise:} \(\tilde{m}_i = m_i + \varepsilon_i\),
  \(\varepsilon_i \sim \mathcal{N}(0, \sigma_{\text{noise}}^2 I)\)
\item
  \textbf{Receive:}
  \(r_i = \frac{1}{K} \sum_{j \in \mathcal{N}(i)} \tilde{m}_j\)
\item
  \textbf{Update:}
  \(s_i' = \sigma\!\big(\alpha \, s_i + (1 - \alpha) \, r_i + \text{drive}\big)\),
  where drive \(\sim \mathcal{N}(0, \delta^2 I)\), \(\delta = 0.02\)
\item
  \textbf{Learn:}
  \(W_i \leftarrow W_i - \eta \, \nabla_{W_i} \frac{1}{K}\sum_{j \in \mathcal{N}(i)} \|m_i - s_j'\|^2\)
\end{enumerate}

State persistence \(\alpha = 0.3\). Learning rate \(\eta = 0.003\). The
learning rule minimizes neighbor prediction error: each agent tries to
make its broadcast message match its neighbors' next states.
Self-prediction, \(\cos(m_i, s_i')\), is never in the gradient.

\subsubsection{2.2 Topologies}\label{topologies}

I test five graph structures:

\begin{itemize}
\tightlist
\item
  \textbf{von Neumann} (\(K = 4\)): cardinal neighbors on a toroidal
  grid
\item
  \textbf{Moore} (\(K = 8\)): cardinal + diagonal neighbors
\item
  \textbf{Hexagonal} (\(K = 6\)): hexagonal lattice
\item
  \textbf{Random} (\(K = 4\)): uniformly sampled random graph
\item
  \textbf{Small-world}: von Neumann + 10\% edge rewiring (Watts \&
  Strogatz, 1998)
\end{itemize}

\subsubsection{2.3 MCH Metrics}\label{mch-metrics}

I measure five quantities per agent per tick:

\begin{itemize}
\tightlist
\item
  \textbf{Self-prediction (self):} \(\cos(m_i, s_i')\), how well the
  broadcast predicts the next state
\item
  \textbf{Integration (\(\Phi\)):} Normalized difference between joint
  and partitioned neighborhood prediction residuals:
  \(\Phi = (\text{parts residual} - \text{joint residual}) / \text{parts residual}\).
  This is a simplified proxy for Tononi's integrated information
  (Tononi, 2004)
\item
  \textbf{Reflexivity (\(R\)):}
  \(\cos(m_i, s_i') - \text{mean}_{j \in \mathcal{N}(i)} \cos(m_i, s_j')\).
  Positive \(R\) means the agent predicts itself better than it predicts
  its neighbors
\item
  \textbf{Temporal persistence (\(T\)):}
  \(T = \text{clip}(1 - \sqrt{\text{Var}_{\text{ema}}(\text{self})}, 0, 1)\).
  How stable self-prediction is over time
\item
  \textbf{Causal efficacy (\(E\)):}
  \(\cos(\Delta s_i, \Delta s_i^{\text{self}})\) where
  \(\Delta s_i^{\text{self}} = \sigma(\alpha s_i + \xi_i) - s_i\) is the
  counterfactual state change without neighbor input
\end{itemize}

\subsubsection{2.4 Activation Functions}\label{activation-functions}

I test three activations to check for universality:

\begin{itemize}
\tightlist
\item
  \textbf{tanh} (default): \(f(x) = \tanh(x)\), range \([-1, 1]\)
\item
  \textbf{Linear:} \(f(x) = \text{clip}(x, -1, 1)\), range \([-1, 1]\)
\item
  \textbf{ReLU:} \(f(x) = \text{clip}(\max(0, x), 0, 1)\), range
  \([0, 1]\)
\end{itemize}

All activation-dependent computations (message transform, state update,
learning gradient) use the configured function and its derivative.

\subsubsection{2.5 Experimental Design}\label{experimental-design}

\begin{longtable}[]{@{}
  >{\raggedright\arraybackslash}p{(\linewidth - 12\tabcolsep) * \real{0.1429}}
  >{\centering\arraybackslash}p{(\linewidth - 12\tabcolsep) * \real{0.1429}}
  >{\centering\arraybackslash}p{(\linewidth - 12\tabcolsep) * \real{0.1429}}
  >{\centering\arraybackslash}p{(\linewidth - 12\tabcolsep) * \real{0.1429}}
  >{\centering\arraybackslash}p{(\linewidth - 12\tabcolsep) * \real{0.1429}}
  >{\centering\arraybackslash}p{(\linewidth - 12\tabcolsep) * \real{0.1429}}
  >{\centering\arraybackslash}p{(\linewidth - 12\tabcolsep) * \real{0.1429}}@{}}
\toprule\noalign{}
\begin{minipage}[b]{\linewidth}\raggedright
Experiment
\end{minipage} & \begin{minipage}[b]{\linewidth}\centering
Runs
\end{minipage} & \begin{minipage}[b]{\linewidth}\centering
Seeds
\end{minipage} & \begin{minipage}[b]{\linewidth}\centering
Sizes
\end{minipage} & \begin{minipage}[b]{\linewidth}\centering
Topologies
\end{minipage} & \begin{minipage}[b]{\linewidth}\centering
Noises
\end{minipage} & \begin{minipage}[b]{\linewidth}\centering
Activation
\end{minipage} \\
\midrule\noalign{}
\endhead
\bottomrule\noalign{}
\endlastfoot
Primary (size 24) & 1,250 & 50 & 24 & all 5 & all 5 & tanh \\
Primary (size 12) & 1,250 & 50 & 12 & all 5 & all 5 & tanh \\
Size 48 + metrics & 250 & 10 & 48 & all 5 & all 5 & tanh \\
Scaling (sizes 12-36) & 1,000 & 10 & 12, 18, 24, 36 & all 5 & all 5 &
tanh \\
Null model (lr = 0) & 250 & 10 & 48 & all 5 & all 5 & tanh \\
ReLU universality & 90 & 10 & 24 & vn, moore, random & 3 & ReLU \\
Linear universality & 90 & 10 & 24 & vn, moore, random & 3 & linear \\
\textbf{Total} & \textbf{4,180} & & & & & \\
\end{longtable}

All runs: 1,000 ticks. Noise levels:
\(\sigma_{\text{noise}} \in \{0.04, 0.08, 0.12, 0.20, 0.30\}\). An
optional Conway's Game of Life (B3/S23) substrate layer provides
structured ``base reality'' input, tested separately at size 24.

\begin{center}\rule{0.5\linewidth}{0.5pt}\end{center}

\subsection{3. Results}\label{results}

\subsubsection{3.1 Structural Effects of
Topology}\label{structural-effects-of-topology}

\textbf{Finding 1: Random graph noise immunity.} On random topologies,
all metrics stay invariant across a 7.5x noise range (\(\sigma = 0.04\)
to \(0.30\)): mean self-prediction varies by \(< 0.003\), \(\Phi\) by
\(< 0.003\), \(R\) by \(0.001\), \(T\) by \(0.0006\). This is explained
by applying the law of large numbers twice: noise cancels over \(K\)
independent neighbors (received signal) and over \(N\) agents
(population mean). Grid topologies do not have this property because
spatially adjacent neighbors are correlated.

\textbf{Finding 2: Noise amplification on high-\(K\) grids.} Moore
topology (\(K = 8\)) shows monotonic increase in self-prediction with
noise: \(+0.018\) at size 24, \(+0.031\) at size 12. The mechanism is
symmetry-breaking: noise decorrelates adjacent agents whose inputs
heavily overlap, creating diversity that enriches the learning gradient.
This needs both noise tolerance (high \(K\)) and redundancy to break
(spatial structure).

\textbf{Finding 3: Small-world \(\Phi\) dissociation.} Small-world
shortcuts boost \(\Phi\) by 13-20\% relative to von Neumann while
leaving self-prediction unchanged (within 0.002). All 50 seeds show
positive \(\Phi\) boost, zero exceptions. Long-range shortcuts introduce
uncorrelated neighbors that make individual (parts) prediction worse
while leaving integrated (joint) prediction intact, which mechanically
increases \(\Phi\).

\subsubsection{3.2 The Autonomy-Predictability
Trade-Off}\label{the-autonomy-predictability-trade-off}

\textbf{Finding 4: Self-determination hurts self-knowledge on high-\(K\)
grids.} On Moore grids (\(K = 8\)), temporal persistence and causal
efficacy predict self-prediction in opposite directions:

\begin{longtable}[]{@{}lccccc@{}}
\toprule\noalign{}
Correlation & Moore & Hex & von Neumann & Random & Small-world \\
\midrule\noalign{}
\endhead
\bottomrule\noalign{}
\endlastfoot
\(r(T, \text{self})\) & \textbf{+0.70} & +0.52 & +0.03 & +0.12 &
+0.22 \\
\(r(E, \text{self})\) & \textbf{-0.71} & -0.71 & -0.12 & -0.14 &
+0.06 \\
\(r(T, E)\) & \textbf{-0.82} & -0.74 & +0.07 & +0.32 & -0.08 \\
\end{longtable}

Agents whose trajectories are externally driven (low \(E\)) but whose
self-models are stable over time (high \(T\)) achieve the best
self-prediction. Here is why: on a Moore grid with 8 neighbors, the
received signal is a smoothed, slowly-varying average. Agents driven by
this signal have predictable futures. Agents driven by their own
nonlinear dynamics are harder for themselves to forecast.
Self-determination is noise from the self-model's perspective.

\textbf{Null model confirmation.} A null model with learning disabled
(\(\eta = 0\)) confirms this is a topological invariant, not a learning
artifact:

\begin{longtable}[]{@{}lccc@{}}
\toprule\noalign{}
Topology & \(r(T,E)\) learning & \(r(T,E)\) null & \(\Delta\) \\
\midrule\noalign{}
\endhead
\bottomrule\noalign{}
\endlastfoot
Moore & -0.82 & -0.79 & 0.03 \\
Hex & -0.74 & -0.69 & 0.05 \\
von Neumann & +0.07 & +0.11 & 0.04 \\
\end{longtable}

The trade-off is geometric: the state update dynamics and topology alone
produce it. Learning at \(\eta = 0.003\) over 1,000 ticks changes
self-prediction by \(< 0.1\%\).

\textbf{Finding 5: Activation universality.} The T-E anti-correlation
survives all three activation functions:

\begin{longtable}[]{@{}lcc@{}}
\toprule\noalign{}
Activation & Moore \(r(T,E)\) & Moore \(r(E, \text{self})\) \\
\midrule\noalign{}
\endhead
\bottomrule\noalign{}
\endlastfoot
tanh & -0.82 & -0.71 \\
Linear & -0.77 & -0.53 \\
ReLU & -0.76 & -0.61 \\
\end{longtable}

The trade-off arises from network structure, not from the specific
nonlinearity. The argument in Section 4 only requires that \(f\) is
bounded.

\subsubsection{\texorpdfstring{3.3 Functional Identity:
\(E \equiv \Phi\) on Structured
Graphs}{3.3 Functional Identity: E \textbackslash equiv \textbackslash Phi on Structured Graphs}}\label{functional-identity-e-equiv-phi-on-structured-graphs}

\textbf{Finding 6:} Causal efficacy and integration are functionally
identical on grids but decouple on random graphs:

\begin{longtable}[]{@{}lccc@{}}
\toprule\noalign{}
Topology & \(r(E, \Phi)\) & Linear slope & RMSE \\
\midrule\noalign{}
\endhead
\bottomrule\noalign{}
\endlastfoot
von Neumann & +0.998 & 0.297 & 0.002 \\
Moore & +0.994 & 0.349 & 0.005 \\
Hex & +0.996 & 0.296 & 0.003 \\
Small-world & +0.984 & 0.193 & 0.004 \\
Random & +0.178 & 0.008 & n/a \\
\end{longtable}

On structured networks, \(\Phi\) measures causal self-determination, not
topology-agnostic information integration. The equivalence breaks on
random graphs because independent neighbor sampling destroys the link
between an agent's self-determination and its neighborhood prediction
structure. This has implications for IIT: \(\Phi\) partitions may
conflate integration with autonomy in any system with spatial
correlations.

\subsubsection{3.4 Dimensionality
Collapse}\label{dimensionality-collapse}

\textbf{Finding 7:} PCA on the five MCH metrics shows that on Moore
grids, the metrics collapse to one dimension:

\begin{longtable}[]{@{}lccc@{}}
\toprule\noalign{}
Activation & PC1 & PC1+PC2 & Dims for 95\% \\
\midrule\noalign{}
\endhead
\bottomrule\noalign{}
\endlastfoot
tanh & \textbf{82.5\%} & 94.8\% & 3 \\
Linear & \textbf{71.8\%} & 89.6\% & 3 \\
ReLU & \textbf{62.1\%} & 82.8\% & 4 \\
\end{longtable}

\noindent\emph{Note:} tanh values are from size-48 Moore grids ($N = 2{,}304$
agents, 250 runs); Linear and ReLU values are from size-24 Moore grids
($N = 576$, 90 runs each). At size 24, tanh PC1 is 67.4\%, lower than
Linear's 71.8\%. The cross-activation comparison therefore reflects both
activation function and system size. Section 3.5 reports uniform-size scaling.

The PC1 loadings form a consistent ``autonomy-predictability axis'':

\begin{longtable}[]{@{}lccc@{}}
\toprule\noalign{}
Metric & tanh & Linear & ReLU \\
\midrule\noalign{}
\endhead
\bottomrule\noalign{}
\endlastfoot
self & -0.37 & -0.27 & -0.35 \\
\(\Phi\) & \textbf{+0.48} & \textbf{+0.52} & \textbf{+0.52} \\
\(R\) & +0.45 & +0.45 & +0.27 \\
\(T\) & \textbf{-0.44} & \textbf{-0.44} & \textbf{-0.48} \\
\(E\) & \textbf{+0.49} & \textbf{+0.52} & \textbf{+0.56} \\
\end{longtable}

\(\{\Phi, R, E\}\) load positive (autonomy pole); \(\{\text{self}, T\}\)
load negative (predictability pole). On random topologies, 4 components
are needed for 95\% variance, so no collapse occurs. The consciousness
correlates are not independent on structured networks. They are
projections of a single underlying variable.

\subsubsection{3.5 Scaling Law}\label{scaling-law}

\textbf{Finding 8 (headline result):} The dimensionality collapse
follows a logarithmic scaling law across five grid sizes:

\begin{longtable}[]{@{}lcccc@{}}
\toprule\noalign{}
Size & \(N\) agents & PC1 (Moore) & PC1 (Random) & \(r(T,E)\) Moore \\
\midrule\noalign{}
\endhead
\bottomrule\noalign{}
\endlastfoot
12 & 144 & 48.8\% & 49.5\% & -0.11 \\
18 & 324 & 55.9\% & 38.7\% & -0.25 \\
24 & 576 & 67.4\% & 36.3\% & -0.70 \\
36 & 1,296 & 75.1\% & 47.3\% & -0.78 \\
48 & 2,304 & 82.5\% & 39.6\% & -0.82 \\
\end{longtable}

The fit is:

\[\text{PC1} = -0.136 + 0.124 \times \ln(N), \quad R^2 = 0.987\]

Random topologies show no scaling (PC1 stays around 40\% at all sizes).
The inflection occurs between \(N = 324\) and \(N = 576\):

\begin{longtable}[]{@{}lcc@{}}
\toprule\noalign{}
Transition & \(\Delta\)PC1 & \(\Delta r(T,E)\) \\
\midrule\noalign{}
\endhead
\bottomrule\noalign{}
\endlastfoot
144 → 324 & +0.071 & -0.14 \\
\textbf{324 → 576} & \textbf{+0.115} & \textbf{-0.45} \\
576 → 1,296 & +0.077 & -0.08 \\
1,296 → 2,304 & +0.074 & -0.04 \\
\end{longtable}

Below about 300 agents, the five metrics are effectively independent
(PC1 around 49\%, \(r(T,E) \approx -0.11\)). Above about 600, they are
locked into a single axis (PC1 \(> 67\%\), \(r(T,E) < -0.70\)). The
constraint on consciousness correlates emerges with scale on structured
topologies. It is a phase-like transition from unconstrained to locked.

\textbf{Full topology comparison across sizes:}

\begin{longtable}[]{@{}lccccc@{}}
\toprule\noalign{}
Topology & PC1 @12 & PC1 @18 & PC1 @24 & PC1 @36 & PC1 @48 \\
\midrule\noalign{}
\endhead
\bottomrule\noalign{}
\endlastfoot
von Neumann & 39.4\% & 46.1\% & 41.5\% & 48.6\% & 45.8\% \\
Moore & 48.8\% & 55.9\% & 67.4\% & 75.1\% & 82.5\% \\
Hex & 47.7\% & 57.9\% & 57.2\% & 71.6\% & 78.4\% \\
Random & 49.5\% & 38.7\% & 36.3\% & 47.3\% & 39.6\% \\
Small-world & 41.3\% & 42.8\% & 41.2\% & 41.3\% & 48.7\% \\
\end{longtable}

Moore and hex diverge sharply from the others as size increases. The
collapse is specific to high-\(K\) structured topologies.

\subsubsection{3.6 Additional Findings}\label{additional-findings}

\textbf{Simpson's paradox.} On Moore and hex, the within-noise
\(\Phi\)-self correlation is positive (\(r \approx +0.48\)) but the
pooled correlation is negative (\(r = -0.28\)). Noise simultaneously
increases self-prediction and decreases \(\Phi\), creating a confound
that reverses the sign. The paradox replicates at sizes 12 and 24 and is
specific to high-connectivity grids.

\textbf{Game of Life substrate.} Adding a Conway's Game of Life
substrate layer (Berlekamp et al., 1982) triples mean self-prediction
(\(0.19 \to 0.57\) at size 24). \(E\) drops (\(0.66 \to 0.44\)),
confirming agents become more externally driven. \(T\) increases
(\(0.66 \to 0.80\)), consistent with the T-E trade-off: structured
external signal leads to lower \(E\), which leads to higher \(T\), which
leads to better self-prediction.

\begin{center}\rule{0.5\linewidth}{0.5pt}\end{center}

\subsection{4. Theoretical Analysis}\label{theoretical-analysis}

\subsubsection{4.1 The Autonomy-Predictability
Argument}\label{the-autonomy-predictability-argument}

\textbf{Claim.} Consider \(N\) agents on a \(K\)-regular graph with
dynamics

\[s_i' = f\!\big(\alpha \, s_i + (1 - \alpha) \, r_i + \xi_i\big)\]

where \(f\) is any bounded activation (\(f: \mathbb{R} \to [-B, B]\)),
\(r_i = \frac{1}{K}\sum_{j \in \mathcal{N}(i)} (g(s_j) + \varepsilon_j)\)
is the received signal, \(\xi_i \sim \mathcal{N}(0, \delta^2 I)\) is
drive noise, and \(\varepsilon_j \sim \mathcal{N}(0, \sigma^2 I)\) is
channel noise. Define causal efficacy
\(E_i = \cos\!\big(s_i' - s_i, \; f(\alpha s_i + \xi_i) - s_i\big)\) and
self-prediction \(\text{self}_i = \cos(m_i, s_i')\).

Then for \(K\) large enough and neighbor correlations \(\rho > 0\),
\(\text{Cov}(E_i, \text{self}_i) < 0\) across the agent population.

\textbf{Proof sketch.}

\emph{Step 1.} Decompose the pre-activation:
\(z_i = \underbrace{\alpha s_i + \xi_i}_{z_i^{\text{self}}} + \underbrace{(1-\alpha) r_i}_{z_i^{\text{ext}}}\).

\emph{Step 2.} On a regular grid with neighbor correlation \(\rho > 0\),
the temporal variance of the external component is:
\[\text{Var}_t(r_i) = \frac{\text{Var}(m)}{K}\big(1 + \rho(K-1)\big) + \frac{\sigma^2}{K}\]
For Moore grids (\(K = 8\), \(\rho \approx 0.3\)):
\(\text{Var}_t(r_i) \approx 0.4 \cdot \text{Var}(m)\).

\emph{Step 3.} The self-driven variance
\(\text{Var}_t(z_i^{\text{self}}) = \alpha^2 \text{Var}_t(s_i) + \delta^2\)
is \(K\)-independent.

\emph{Step 4.} For large \(K\) with \(\rho > 0\),
\(\text{Var}(z^{\text{ext}}) < \text{Var}(z^{\text{self}})\). Agents
dominated by the external component (low \(E\)) have smoother
trajectories. Smoother trajectories are easier to self-predict.
Therefore \(\text{Cov}(E, \text{self}) < 0\).

\emph{Step 5.} On random graphs (\(\rho = 0\)), the external signal is
just averaged noise from independent sources, which is unreliable rather
than predictable. The negative covariance does not emerge.

\emph{Step 6.} The argument uses only: (i) \(f\) is bounded, (ii)
\(r_i\) is an average of \(K\) messages, (iii) \(\rho > 0\). No property
specific to tanh is needed.

\subsubsection{4.2 Random Graph Fixed
Point}\label{random-graph-fixed-point}

Random graphs are immune to noise because of a self-stabilizing loop. On
random graphs, each agent's \(K\) neighbors are independent uniform
samples from the population. The received signal
\(r_i \approx \bar{m} + \mathcal{N}(0, \sigma^2/K \cdot I)\) converges
to the population mean \(\bar{m}\) by the law of large numbers. As noise
\(\sigma\) increases, the zero-mean perturbations cancel across both the
\(K\) neighbors (received signal) and \(N\) agents (population mean).
The expected state, gradient direction, and all macroscopic statistics
stay the same (see Appendix for full derivation).

\subsubsection{\texorpdfstring{4.3 Small-World \(\Phi\)
Dissociation}{4.3 Small-World \textbackslash Phi Dissociation}}\label{small-world-phi-dissociation}

Long-range shortcuts connect agents to distant, uncorrelated partners.
Each shortcut makes the parts prediction worse (the distant agent is a
poor individual predictor) while leaving the joint prediction intact
(the mean is dominated by local neighbors). Since
\(\Phi \propto (\text{parts} - \text{joint}) / \text{parts}\), shortcuts
mechanically increase \(\Phi\) without changing self-prediction (the
long-range message is about 10\% of the input and zero-mean in
direction).

\begin{center}\rule{0.5\linewidth}{0.5pt}\end{center}

\subsection{5. Discussion}\label{discussion}

\subsubsection{5.1 The Autonomy-Predictability
Principle}\label{the-autonomy-predictability-principle}

The central finding of this work flips a deep assumption in the
consciousness literature. IIT (Tononi, 2004; Tononi et al., 2016), the
MCH (Fitz, 2025), attention schema theory (Graziano, 2013), and most
autonomy-based frameworks assume that self-determination enables
self-knowledge. My data show the opposite on high-connectivity networks:
the agents most in control of their own trajectories are the worst at
predicting themselves (\(r(E, \text{self}) = -0.71\)). The reason is
simple: external determination is regularizing. When an agent's future
state is driven by the average of 8 predictable neighbors, that future
is easy to forecast. When it is driven by the agent's own nonlinear
dynamics, it is chaotic and hard to predict.

The activation universality (Finding 5), null model confirmation, and
formal argument (Section 4.1) together give strong evidence that this is
a structural property of networked systems, not an artifact of any
particular implementation choice.

\subsubsection{5.2 Implications for Integrated Information
Theory}\label{implications-for-integrated-information-theory}

The \(E \equiv \Phi\) identity (Finding 6) raises questions for IIT that
are worth investigating. On structured networks in my system, the
\(\Phi\) proxy (which was designed to approximate information
integration) tracks causal self-determination almost perfectly
(\(r > 0.994\)). If this relationship holds for more rigorous \(\Phi\)
measures, it would suggest that on systems with spatially correlated
neighbors, \(\Phi\) partitions may conflate integration with autonomy.
The decoupling on random graphs (\(r = 0.18\)) shows this is not a
universal property but a consequence of structured connectivity. I want
to be clear that this is only demonstrated for my \(\Phi\) proxy (see
Section 5.6), and testing with IIT-compliant \(\Phi\) measures is an
important next step.

\subsubsection{5.3 Dimensionality Collapse and
Overcounting}\label{dimensionality-collapse-and-overcounting}

The dimensionality collapse (Finding 7) shows that on structured
networks, the four MCH correlates (\(\Phi\), \(R\), \(T\), \(E\)) are
not independent indicators of consciousness. They are projections of a
single underlying variable: the agent's position on the
autonomy-predictability axis. A theory that treats them as separate
lines of evidence overcounts the degrees of freedom. Adding more
correlates does not add more information. They measure the same thing
through different lenses.

This result depends on topology: on random networks, where neighbor
signals are independent, the five metrics keep their independence (4
components for 95\% variance). The constraint is not universal. It is an
emergent property of structured connectivity.

\subsubsection{5.4 The Scaling Law as Phase
Transition}\label{the-scaling-law-as-phase-transition}

The scaling law (Finding 8) answers the question that motivated this
project: does anything \emph{emerge} from scale? The answer is that the
\textbf{constraint} on consciousness correlates emerges from scale.
Small structured systems (\(N < 300\)) are unconstrained, with five
metrics and five degrees of freedom. Large ones (\(N > 600\)) are forced
into a single trade-off axis. The transition sharpens between
\(N = 324\) and \(N = 576\), with the largest gains in both PC1
(\(+0.115\)) and \(r(T,E)\) (\(-0.45\)) happening in this interval.

The logarithmic fit (\(R^2 = 0.987\)) predicts PC1 \(\to 1.0\) (complete
collapse) at \(N \approx 8{,}000\) agents, though the curve has to
saturate before that point. Random systems never undergo this
transition, confirming that both structure and scale are necessary.

This is analogous to phase transitions in statistical physics: below a
critical system size, macroscopic constraints do not show up. Above it,
microscopic degrees of freedom become coupled by the network geometry.
Khajehabdollahi (2018) found that \(\Phi\) undergoes a phase transition
at critical temperature in Ising models, which is interesting to compare
with. But my result is different in kind: I show the emergence of
constraint \emph{across} correlates, not a transition of a single
measure.

\subsubsection{5.5 Related Work}\label{related-work}

\textbf{Integrated Information Theory.} Tononi's IIT (Tononi, 2004;
Tononi et al., 2016; Albantakis et al., 2023) proposes \(\Phi\) as the
primary measure of consciousness. Computational studies have explored
\(\Phi\) in Ising models (Khajehabdollahi, 2018), coupled repressilators
(Abrego \& Zaikin, 2019), fish schools (Niizato et al., 2024), and
neuropercolation models (Luitle, 2023). These works study \(\Phi\) in
isolation. I measure \(\Phi\) alongside four other correlates and show
it is not independent of them on structured networks.

\textbf{Predictive processing.} The free energy principle (Friston,
2010) and predictive coding (Rao \& Ballard, 1999; Clark, 2013) propose
that neural systems minimize prediction error. My agents do a
distributed version of this: each one minimizes neighbor prediction
error. The emergent self-prediction is consistent with predictive
processing accounts of self-awareness (Seth, 2014; Seth \& Bayne, 2022;
Hohwy, 2013).

\textbf{Network topology and dynamics.} The importance of topology for
collective dynamics is well established (Watts \& Strogatz, 1998;
Barabasi \& Albert, 1999; Newman, 2003). I extend this to
consciousness-relevant metrics, showing that topology controls not just
individual metrics but their joint degrees of freedom.

\textbf{Multi-agent emergence.} Arsiwalla et al.~(2023) propose a
``morphospace of consciousness'' mapping complexity dimensions. Bailey
and Schneider (2025) propose spectral measures of epistemic emergence
with topological constraints. Both are theoretical. My work provides the
first simulation-based evidence that consciousness correlates collapse
to a low-dimensional manifold on structured networks.

\textbf{Machine consciousness.} Fitz (2025) proposes the MCH as a
theoretical program. My work is a partial computational implementation.
I borrow the MCH's metric definitions and communication framework but I
lack the second-order perception Fitz considers essential. I find that
one of the MCH's implicit assumptions, that \(\Phi\), \(R\), \(T\),
\(E\) provide independent evidence, does not hold on the structured
networks I test.

\subsubsection{5.6 Limitations}\label{limitations}

Several limitations constrain the strength of my claims. I start with
the one most likely to concern reviewers.

\textbf{\(R\) shares variance with self-prediction by construction.}
Reflexivity is defined as
\(R = \cos(m_i, s_i') - \text{mean}_{j} \cos(m_i, s_j')\), so \(R\) and
self-prediction share the \(\cos(m_i, s_i')\) term. If this constructive
overlap drove the PCA collapse, the headline finding would be partly an
artifact of metric definitions rather than a property of the system. I
tested this directly by repeating all PCA analyses with \(R\) excluded.
On Moore grids at size 48, PC1 \emph{without} \(R\) is 83.9\%, compared
to 82.5\% \emph{with} \(R\). Including \(R\) does not inflate the
collapse. Across all five scaling sizes, removing \(R\) changes PC1 by
less than 2 percentage points in either direction (size 12: 49.1\% vs
48.8\%; size 24: 69.8\% vs 67.4\%; size 48: 83.9\% vs 82.5\%). The
dimensionality collapse is robust to this exclusion. I report results
with all five metrics because \(R\) carries interpretive value
(self-vs-neighbor prediction asymmetry), but the collapse holds on four
constructively independent metrics. The structural concern remains:
these are not five fully independent measurements, and I do not claim
they are.

\textbf{The \(\Phi\) proxy is not IIT's \(\Phi\).} My integration
measure, the normalized gap between joint and partitioned prediction
residuals, is a simplified heuristic, not Tononi's integrated
information (computing exact \(\Phi\) is NP-hard; Tegmark, 2016). The
\(E \equiv \Phi\) identity (Finding 6) may therefore reflect a
constructive relationship between my proxy and my efficacy measure,
rather than a deep property of Tononi's \(\Phi\). My implications for
IIT (Section 5.2) should be read as hypotheses that motivate testing
with more rigorous \(\Phi\) approximations, not as established results
about IIT.

\textbf{The theorem is a proof sketch.} The argument in Section 4.1
establishes the intuitive mechanism (neighbor-averaging smooths
trajectories, making externally-driven agents more self-predictable) but
Step 4 (``smoother trajectories are easier to self-predict'') is an
empirically supported assertion, not a formally derived inequality. A
rigorous version would need to bound \(\text{Cov}(E, \text{self})\) in
terms of \(K\) and \(\rho\), which I have not done. I call it a
``theorem'' in the informal mathematical sense of a structured argument,
not a complete formal proof.

\textbf{The scaling law is fit to five data points.} Any smooth
monotonic function can fit a logarithm with high \(R^2\) over a single
decade (\(N = 144\) to \(2{,}304\)). I cannot tell a log apart from a
power law or sigmoid on this range. The \(R^2 = 0.987\) shows good fit
to the data I have, not confirmation of the functional form. The
extrapolation to \(N \approx 8{,}000\) (Section 5.4) is speculative.

\textbf{Three activations is limited universality.} I tested tanh,
clipped-linear, and clipped-ReLU. Two share the output range
\([-1, 1]\). The ReLU and linear sweeps use 90 runs each on 3
topologies, which is far less data than the primary tanh sweeps. I have
not tested sigmoid, GELU, or unbounded activations. The universality
claim is that the trade-off does not require tanh specifically, not that
it survives all possible nonlinearities.

\textbf{My agents are simple.} Each agent has 64 parameters
(\(8 \times 8\) weight matrix) and 8-dimensional state. This is nowhere
near comparable to biological neural circuits or transformer-based
architectures. My claims about ``consciousness correlates'' apply within
this system. They are structural properties of networked agents with
these dynamics. Extrapolation to biological or sophisticated artificial
systems is speculative and should be treated as a hypothesis for future
work, not an established finding.

\textbf{consim is not a full MCH implementation.} Fitz's MCH envisions
transformer-based agents capable of second-order perception. My agents
only do first-order prediction (predicting neighbors, not predicting
their own predictions). I borrow the MCH's metric definitions and
communication framework, but the system lacks the representational
capacity the MCH considers necessary for consciousness.

\textbf{Other limitations.} All grids use toroidal (periodic) boundary
conditions, which eliminates edge effects present in real networks. The
Game of Life substrate coupling is additive rather than perceptual. The
scaling-sweep experiments use 10 seeds (vs.~50 for the primary sweeps),
providing less statistical power for the scaling law. I do not claim
consciousness emergence. My system is too simple for that.

What I do show is that the \emph{structural logic} of self-prediction
from predictive communication produces non-trivial, topology-dependent
phenomena that constrain which combinations of consciousness-relevant
metrics are jointly achievable, and that these constraints scale with
system size on structured networks.

\begin{center}\rule{0.5\linewidth}{0.5pt}\end{center}

\subsection{6. Conclusion}\label{conclusion}

I showed that on structured networks, five consciousness-relevant
metrics (self-prediction, \(\Phi\), \(R\), \(T\), and \(E\)) are not
independent. They collapse to a single ``autonomy-predictability axis''
whose strength follows a logarithmic scaling law with system size. This
collapse works like an uncertainty principle for networked systems: on
\(K\)-regular graphs with correlated neighbors, an agent cannot
simultaneously maximize causal autonomy and self-predictability. The
constraint emerges above a critical system size (\(N \approx 300\) to
\(600\) on Moore grids) and is absent on random topologies at any size.

Three implications follow:

\begin{enumerate}
\def\labelenumi{\arabic{enumi}.}
\item
  \textbf{For consciousness theory:} Any framework that treats
  integration, reflexivity, persistence, efficacy, and self-knowledge as
  independent correlates overcounts degrees of freedom on structured
  networks. A single number, the agent's position on the
  autonomy-predictability axis, captures 62--83\% of all variation
  (across activations and system sizes).
\item
  \textbf{For IIT (tentative):} On structured networks, my \(\Phi\)
  proxy and causal efficacy are functionally identical (\(r > 0.994\)).
  If this extends to rigorous \(\Phi\) measures, it would suggest
  \(\Phi\) conflates self-determination with information integration
  when neighbors are correlated.
\item
  \textbf{For the MCH:} In a partial implementation of Fitz's framework
  (see Section 5.6), the four correlates cannot be simultaneously
  maximized. They lie on a Pareto frontier determined by network
  topology.
\end{enumerate}

Code, data (4,180 runs), and all analysis scripts are publicly available
at github.com/brian-mwirigi/consim (code under MIT license).

\textbf{Acknowledgment.} AI tools (GitHub Copilot) assisted with coding.

\subsection*{Errata (v2)}

The following corrections were made relative to v1 (DOI: 10.5281/zenodo.18940956):

\begin{enumerate}
\item \textbf{Section 5.6, \(R\)-exclusion PCA (6 numbers corrected).}
  v1 reported stale values. Corrected:
  size~12 with/without~\(R\): 48.8\%/49.1\% (was 49.9\%/49.3\%);
  size~24: 67.4\%/69.8\% (was 66.7\%/68.7\%);
  size~48: 82.5\%/83.9\% (was 85.0\%/85.7\%).
  Qualitative conclusion unchanged.
\item \textbf{Section 3.4, size mismatch disclosed.}
  The Finding~7 PCA table compared tanh at size~48 with Linear/ReLU at
  size~24 without disclosure. A note now states that at the same size~24,
  tanh PC1 is 67.4\%, lower than Linear's 71.8\%.
\item \textbf{Section 6, ``62--83\%'' clarified.}
  Range spans different activations \emph{and} system sizes; parenthetical
  added.
\end{enumerate}

\begin{center}\rule{0.5\linewidth}{0.5pt}\end{center}

\subsection{References}\label{references}

Abrego, L., \& Zaikin, A. (2019). Integrated information as a measure of
cognitive processes in coupled genetic repressilators. \emph{Entropy},
21(4), 382.

Albantakis, L., Barbosa, L., Findlay, G., Grasso, M., Haun, A. M.,
Marshall, W., \ldots{} \& Tononi, G. (2023). Integrated information
theory (IIT) 4.0: formulating the properties of phenomenal existence in
physical terms. \emph{PLoS Computational Biology}, 19(10), e1011465.

Arsiwalla, X. D., Sole, R., Moulin-Frier, C., \& Herreros, I. (2023).
The morphospace of consciousness: Three kinds of complexity for minds
and machines. \emph{NeuroSci}, 4(2), 79-102.

Baars, B. J. (1988). \emph{A Cognitive Theory of Consciousness}.
Cambridge University Press.

Bailey, M., \& Schneider, S. (2025). When wholes resist decomposition: A
spectral measure of epistemic emergence. \emph{Preprint}.

Barabasi, A.-L., \& Albert, R. (1999). Emergence of scaling in random
networks. \emph{Science}, 286(5439), 509-512.

Berlekamp, E. R., Conway, J. H., \& Guy, R. K. (1982). \emph{Winning
Ways for Your Mathematical Plays}. Academic Press.

Clark, A. (2013). Whatever next? Predictive brains, situated agents, and
the future of cognitive science. \emph{Behavioral and Brain Sciences},
36(3), 181-204.

Cleeremans, A. (2011). The radical plasticity thesis: how the brain
learns to be conscious. \emph{Frontiers in Psychology}, 2, 86.

Fitz, S. (2025). Testing the Machine Consciousness Hypothesis.
\emph{arXiv:2512.01081}.

Friston, K. (2010). The free-energy principle: a unified brain theory?
\emph{Nature Reviews Neuroscience}, 11(2), 127-138.

Graziano, M. S. A. (2013). \emph{Consciousness and the Social Brain}.
Oxford University Press.

Hohwy, J. (2013). \emph{The Predictive Mind}. Oxford University Press.

Khajehabdollahi, S. (2018). Phase transitions of integrated information
in the generalized Ising model of the brain. \emph{Master's thesis,
Western University}.

Luitle, T. (2023). Synchronisation through integrated information in
neuropercolation models near criticality. \emph{Master's thesis}.

Newman, M. E. J. (2003). The structure and function of complex networks.
\emph{SIAM Review}, 45(2), 167-256.

Niizato, T., Sakamoto, K., Mototake, Y., Murakami, H., Tomaru, T.,
Hoshika, Y., \& Fukushima, T. (2024). Information structure of
heterogeneous criticality in a fish school. \emph{Scientific Reports},
14, 23503.

Rao, R. P. N., \& Ballard, D. H. (1999). Predictive coding in the visual
cortex: a functional interpretation of some extra-classical
receptive-field effects. \emph{Nature Neuroscience}, 2(1), 79-87.

Seth, A. K. (2014). A predictive processing theory of sensorimotor
contingencies: Explaining the puzzle of perceptual presence and its
absence in synesthesia. \emph{Cognitive Neuroscience}, 5(2), 97-118.

Seth, A. K., \& Bayne, T. (2022). Theories of consciousness.
\emph{Nature Reviews Neuroscience}, 23, 439-452.

Tegmark, M. (2016). Improved measures of integrated information.
\emph{PLoS Computational Biology}, 12(11), e1005123.

Tononi, G. (2004). An information integration theory of consciousness.
\emph{BMC Neuroscience}, 5(42).

Tononi, G., Boly, M., Massimini, M., \& Koch, C. (2016). Integrated
information theory: from consciousness to its physical substrate.
\emph{Nature Reviews Neuroscience}, 17(7), 450-461.

Watts, D. J., \& Strogatz, S. H. (1998). Collective dynamics of
`small-world' networks. \emph{Nature}, 393(6684), 440-442.

\begin{center}\rule{0.5\linewidth}{0.5pt}\end{center}

\subsection{Appendix A:
Reproducibility}\label{appendix-a-reproducibility}

All experiments can be reproduced with the following commands:

\begin{Shaded}
\begin{Highlighting}[]
\ExtensionTok{pip}\NormalTok{ install numpy matplotlib Pillow}

\CommentTok{\# Primary sweeps (sizes 12 and 24)}
\ExtensionTok{python}\NormalTok{ run.py }\AttributeTok{{-}{-}sweep} \AttributeTok{{-}{-}sweep{-}seeds}\NormalTok{ 1{-}50 }\DataTypeTok{\textbackslash{}}
  \AttributeTok{{-}{-}sweep{-}topos}\NormalTok{ von\_neumann,moore,hex,random,small\_world }\DataTypeTok{\textbackslash{}}
  \AttributeTok{{-}{-}sweep{-}noises}\NormalTok{ 0.04,0.08,0.12,0.20,0.30 }\DataTypeTok{\textbackslash{}}
  \AttributeTok{{-}{-}ticks}\NormalTok{ 1000 }\AttributeTok{{-}{-}size}\NormalTok{ 24 }\AttributeTok{{-}{-}sweep{-}csv}\NormalTok{ sweep\_1250.csv}

\ExtensionTok{python}\NormalTok{ run.py }\AttributeTok{{-}{-}sweep} \AttributeTok{{-}{-}sweep{-}seeds}\NormalTok{ 1{-}50 }\DataTypeTok{\textbackslash{}}
  \AttributeTok{{-}{-}sweep{-}topos}\NormalTok{ von\_neumann,moore,hex,random,small\_world }\DataTypeTok{\textbackslash{}}
  \AttributeTok{{-}{-}sweep{-}noises}\NormalTok{ 0.04,0.08,0.12,0.20,0.30 }\DataTypeTok{\textbackslash{}}
  \AttributeTok{{-}{-}ticks}\NormalTok{ 1000 }\AttributeTok{{-}{-}size}\NormalTok{ 12 }\AttributeTok{{-}{-}sweep{-}csv}\NormalTok{ sweep\_size12.csv}

\CommentTok{\# Size 48 with R, T, E metrics}
\ExtensionTok{python}\NormalTok{ run.py }\AttributeTok{{-}{-}sweep} \AttributeTok{{-}{-}sweep{-}seeds}\NormalTok{ 1{-}10 }\DataTypeTok{\textbackslash{}}
  \AttributeTok{{-}{-}sweep{-}topos}\NormalTok{ von\_neumann,moore,hex,random,small\_world }\DataTypeTok{\textbackslash{}}
  \AttributeTok{{-}{-}sweep{-}noises}\NormalTok{ 0.04,0.08,0.12,0.20,0.30 }\DataTypeTok{\textbackslash{}}
  \AttributeTok{{-}{-}ticks}\NormalTok{ 1000 }\AttributeTok{{-}{-}size}\NormalTok{ 48 }\AttributeTok{{-}{-}sweep{-}csv}\NormalTok{ sweep\_size48.csv}

\CommentTok{\# Null model (learning disabled)}
\ExtensionTok{python}\NormalTok{ run.py }\AttributeTok{{-}{-}sweep} \AttributeTok{{-}{-}sweep{-}seeds}\NormalTok{ 1{-}10 }\DataTypeTok{\textbackslash{}}
  \AttributeTok{{-}{-}sweep{-}topos}\NormalTok{ von\_neumann,moore,hex,random,small\_world }\DataTypeTok{\textbackslash{}}
  \AttributeTok{{-}{-}sweep{-}noises}\NormalTok{ 0.04,0.08,0.12,0.20,0.30 }\DataTypeTok{\textbackslash{}}
  \AttributeTok{{-}{-}ticks}\NormalTok{ 1000 }\AttributeTok{{-}{-}size}\NormalTok{ 48 }\AttributeTok{{-}{-}lr}\NormalTok{ 0 }\AttributeTok{{-}{-}sweep{-}csv}\NormalTok{ sweep\_null.csv}

\CommentTok{\# Scaling sweeps}
\ExtensionTok{python}\NormalTok{ run.py }\AttributeTok{{-}{-}sweep} \AttributeTok{{-}{-}sweep{-}seeds}\NormalTok{ 1{-}10 }\DataTypeTok{\textbackslash{}}
  \AttributeTok{{-}{-}sweep{-}topos}\NormalTok{ von\_neumann,moore,hex,random,small\_world }\DataTypeTok{\textbackslash{}}
  \AttributeTok{{-}{-}sweep{-}noises}\NormalTok{ 0.04,0.08,0.12,0.20,0.30 }\DataTypeTok{\textbackslash{}}
  \AttributeTok{{-}{-}ticks}\NormalTok{ 1000 }\AttributeTok{{-}{-}size}\NormalTok{ 12 }\AttributeTok{{-}{-}sweep{-}csv}\NormalTok{ sweep\_size12\_full.csv}

\ExtensionTok{python}\NormalTok{ run.py }\AttributeTok{{-}{-}sweep} \AttributeTok{{-}{-}sweep{-}seeds}\NormalTok{ 1{-}10 }\DataTypeTok{\textbackslash{}}
  \AttributeTok{{-}{-}sweep{-}topos}\NormalTok{ von\_neumann,moore,hex,random,small\_world }\DataTypeTok{\textbackslash{}}
  \AttributeTok{{-}{-}sweep{-}noises}\NormalTok{ 0.04,0.08,0.12,0.20,0.30 }\DataTypeTok{\textbackslash{}}
  \AttributeTok{{-}{-}ticks}\NormalTok{ 1000 }\AttributeTok{{-}{-}size}\NormalTok{ 18 }\AttributeTok{{-}{-}sweep{-}csv}\NormalTok{ sweep\_size18\_full.csv}

\ExtensionTok{python}\NormalTok{ run.py }\AttributeTok{{-}{-}sweep} \AttributeTok{{-}{-}sweep{-}seeds}\NormalTok{ 1{-}10 }\DataTypeTok{\textbackslash{}}
  \AttributeTok{{-}{-}sweep{-}topos}\NormalTok{ von\_neumann,moore,hex,random,small\_world }\DataTypeTok{\textbackslash{}}
  \AttributeTok{{-}{-}sweep{-}noises}\NormalTok{ 0.04,0.08,0.12,0.20,0.30 }\DataTypeTok{\textbackslash{}}
  \AttributeTok{{-}{-}ticks}\NormalTok{ 1000 }\AttributeTok{{-}{-}size}\NormalTok{ 24 }\AttributeTok{{-}{-}sweep{-}csv}\NormalTok{ sweep\_size24\_full.csv}

\ExtensionTok{python}\NormalTok{ run.py }\AttributeTok{{-}{-}sweep} \AttributeTok{{-}{-}sweep{-}seeds}\NormalTok{ 1{-}10 }\DataTypeTok{\textbackslash{}}
  \AttributeTok{{-}{-}sweep{-}topos}\NormalTok{ von\_neumann,moore,hex,random,small\_world }\DataTypeTok{\textbackslash{}}
  \AttributeTok{{-}{-}sweep{-}noises}\NormalTok{ 0.04,0.08,0.12,0.20,0.30 }\DataTypeTok{\textbackslash{}}
  \AttributeTok{{-}{-}ticks}\NormalTok{ 1000 }\AttributeTok{{-}{-}size}\NormalTok{ 36 }\AttributeTok{{-}{-}sweep{-}csv}\NormalTok{ sweep\_size36\_full.csv}

\CommentTok{\# Activation universality}
\ExtensionTok{python}\NormalTok{ run.py }\AttributeTok{{-}{-}sweep} \AttributeTok{{-}{-}sweep{-}seeds}\NormalTok{ 1{-}10 }\DataTypeTok{\textbackslash{}}
  \AttributeTok{{-}{-}sweep{-}topos}\NormalTok{ von\_neumann,moore,random }\DataTypeTok{\textbackslash{}}
  \AttributeTok{{-}{-}sweep{-}noises}\NormalTok{ 0.04,0.12,0.30 }\DataTypeTok{\textbackslash{}}
  \AttributeTok{{-}{-}ticks}\NormalTok{ 1000 }\AttributeTok{{-}{-}size}\NormalTok{ 24 }\AttributeTok{{-}{-}activation}\NormalTok{ relu }\AttributeTok{{-}{-}sweep{-}csv}\NormalTok{ sweep\_relu.csv}

\ExtensionTok{python}\NormalTok{ run.py }\AttributeTok{{-}{-}sweep} \AttributeTok{{-}{-}sweep{-}seeds}\NormalTok{ 1{-}10 }\DataTypeTok{\textbackslash{}}
  \AttributeTok{{-}{-}sweep{-}topos}\NormalTok{ von\_neumann,moore,random }\DataTypeTok{\textbackslash{}}
  \AttributeTok{{-}{-}sweep{-}noises}\NormalTok{ 0.04,0.12,0.30 }\DataTypeTok{\textbackslash{}}
  \AttributeTok{{-}{-}ticks}\NormalTok{ 1000 }\AttributeTok{{-}{-}size}\NormalTok{ 24 }\AttributeTok{{-}{-}activation}\NormalTok{ linear }\AttributeTok{{-}{-}sweep{-}csv}\NormalTok{ sweep\_linear.csv}
\end{Highlighting}
\end{Shaded}

Raw data (10 CSV files, 4,180 runs total) available at
github.com/brian-mwirigi/consim.

\end{document}
